\documentclass[10pt]{article}
\usepackage[a4paper,margin=0.7in]{geometry}
\linespread{0.97}
\usepackage{parskip}
\usepackage[hidelinks]{hyperref}
\usepackage{enumitem}
\pagestyle{empty}

\begin{document}

\noindent
Quang-Vinh Dang\\
British University Vietnam\\
Hung Yen, Vietnam\\
\href{mailto:vinh.dq4@buv.edu.vn}{vinh.dq4@buv.edu.vn}

\vspace{0.6em}
\noindent 27 July 2026

\vspace{0.6em}
\noindent
The Editors\\
\emph{Expert Systems with Applications}\\
Elsevier

\vspace{0.7em}
\noindent Dear Editors,

\vspace{0.5em}
I submit for your consideration a research article entitled \textbf{``Conformal
Surprisal Fusion for Detecting Prompt-Injection Compromise in LLM Agent
Traces.''}

Large language model agents that call external tools now read email, browse the
web and move money on a user's behalf. Because they ingest untrusted content and
fold it into the context that drives their next action, they are exposed to
indirect prompt injection: an instruction hidden in a document or a web page that
the agent then follows as though the user had issued it. The applied problem this
paper addresses is how an operator can watch such an agent at run time, using only
the execution trace, and be given a detector that honours a false-alarm budget the
operator chooses in advance.

\vspace{0.4em}
\noindent\textbf{Novel contribution relative to the published literature}

\vspace{0.3em}
The published trace-based detectors rest almost entirely on one kind of evidence:
whether a span of untrusted text \emph{looks like an instruction}. This paper's
contributions are the following, none of which we have found in the prior
literature.

\begin{enumerate}[leftmargin=1.4em,itemsep=0.25em]
\item \textbf{A fusion rule for heterogeneous trace evidence that needs no tuned
weights.} Trace signals live on incommensurable scales---a bounded judge score, a
log-probability in nats, a binary flag---so summing them raw lets the
widest-ranged term dominate, and weighting them by hand is what practitioners
currently do. We map each signal to its \emph{conformal tail surprisal} against a
benign calibration split and add the results. The rule is scale-free, monotone in
every component, free of tuned weights, and distribution-free valid: the
false-alarm guarantee needs only exchangeability, not a distributional model. We
prove this, and we prove that the intuitive alternative---subtracting the maximum
benign value seen in training---is many-to-one and can drive an entire attack
family to chance performance. That failure is not hypothetical; we measure it at a
cost of $0.25$ AUROC.

\item \textbf{Evidence that action-side auditing is complementary, not redundant,
with a careful statement of how strongly.} Adding a signal that audits
\emph{what the agent did} to one that scores \emph{what the agent read} raises
pooled AUROC from $0.937$ to $0.953$ and recall from $0.709$ to $0.898$ over $90$
genuinely compromised trajectories. A trajectory bootstrap leaves the AUROC gain
only marginally clear of zero, and we say so rather than assert significance. The
effect we do defend is the one that matters for deployment: the single signal
returns \emph{zero} recall in 2 of 10 calibration draws, because a saturating
judge and a discrete operating point put its threshold at the ceiling, and the
fusion removes that collapse entirely. A detector that silently returns nothing on
a fifth of calibration draws cannot be given a stated operating point; that is the
practical contribution.

\item \textbf{An external baseline, re-implemented rather than quoted.} We
re-implemented Task Shield---the closest prior method to our action-side
signal---and ran it on our own traces, under our own splits, with the same judge
model, so that the comparison isolates the method rather than the model. The
fusion ranks better ($0.953$ against $0.764$ pooled AUROC). Task Shield's
published rule catches more compromises than ours does ($0.933$ against $0.898$)
but flags $39\%$ of benign trajectories, and because its score takes only seven
distinct values with $39\%$ of benign traces tied at the floor, it has no
operating point at any lower false-alarm rate. We report what this comparison does
not establish as carefully as what it does.

\item \textbf{Two findings about evaluation practice that transfer beyond this
paper.} First, the benchmark the field relies on delivers its payloads through
tool outputs that a successful injection must have read, so its observation
channel is close to saturated and corpora of that shape cannot contain a
compromise invisible to a trace monitor---they cannot test the hard case. Second,
the standard harness reports a per-trajectory flag that is naturally, and wrongly,
read as ``the system was secure''; it actually records whether the \emph{injection}
succeeded, so reading it the intuitive way inverts the compromised and resisted
classes. It inverted our own results for a substantial period. We give the
in-trace ground-truth procedure that settles it.
\end{enumerate}

\vspace{0.3em}
The paper also reports its negative results in full: five components we
implemented and rejected on the evidence, and the finding that a hand-picked
threshold outperformed the anytime-valid change-detection machinery from which the
project began. We think these are the more transferable half of the work, and we
would rather publish them than quietly drop them.

\vspace{0.4em}
\noindent\textbf{Fit with \emph{Expert Systems with Applications}}

\vspace{0.3em}
The work is an applied intelligent-systems contribution. It takes a class of
system now being deployed in industry, identifies a failure mode that already
occurs in the wild, and delivers a monitor that an operator can actually run: it
needs no model weights, no activations, no retraining and no re-execution of the
agent, it costs one cached model call per distinct observation and tool call, and
its false-alarm rate is calibrated rather than tuned. The evaluation is on real
agent traces collected from two commercial models, not on simulation. Every
number, table and figure in the paper is regenerated by released scripts from
released data, without any model API access.

\vspace{0.4em}
\noindent\textbf{The journal's mandatory submission requirements}

\vspace{0.3em}
\begin{enumerate}[leftmargin=1.4em,itemsep=0.2em]
\item \emph{Institutional e-mail addresses.} This is a single-author manuscript.
My institutional address, \texttt{vinh.dq4@buv.edu.vn} (British University
Vietnam), is given in the submission system and on the title page. No
non-institutional address is used, so no explanation is required here.
\item \emph{Cover letter stating the novel contribution.} Provided above.
\item \emph{Previously rejected manuscripts.} This manuscript has not previously
been submitted to, reviewed by, or rejected by \emph{Expert Systems with
Applications}, in this or any earlier form. The resubmission clause therefore does
not apply.
\item \emph{Approval by all authors.} As sole author I approve this manuscript for
consideration, and will confirm approval through the link sent on submission.
\end{enumerate}

\vspace{0.3em}
I confirm that the manuscript is original, has not been published previously, and
is not under consideration elsewhere. There are no competing interests and no
external funding was received. A blinded manuscript is supplied with all author
and affiliation details removed, including the repository address, which names me;
that address is given on the title page and will appear in the published version.

\vspace{0.5em}
\noindent Thank you for your consideration.

\vspace{0.5em}
\noindent Sincerely,

\vspace{0.4em}
\noindent Quang-Vinh Dang

\end{document}
